\documentclass[12pt]{article}
\usepackage[utf8]{inputenc}
\usepackage[T1]{fontenc}
\usepackage{hyperref}
\title{Ere Language Resources for K-12 Curriculum Development}
\author{Compiled by Claude Sonnet 5 (claude-sonnet-5) for Andr\'as Kornai}
\date{}
\begin{document}
\maketitle

\section{Introduction}
Ere (also called Nane; ISO 639-3 twp, Glottolog eree1241) is an Austronesian (Oceanic, Manus) language spoken on the south coast of Manus Island, Papua New Guinea. Wikipedia lists about 1,400 speakers, citing the 2000 census, and 1,030 for 1980. The Manus languages form a dialect continuum, so language boundaries are fuzzy. Little material turned up in search.

\section{Primary Educational Resources}
None found.

\section{Linguistic and Reference Materials}
\begin{itemize}
\item Ere language, Wikipedia: \url{https://en.wikipedia.org/wiki/Ere_language}
\item Manus languages, Wikipedia: \url{https://en.wikipedia.org/wiki/Manus_languages}
\item Numerals (Blust data), Max Planck EVA database: \url{https://lingweb.eva.mpg.de/channumerals/Ere.htm}
\item Wiktionary, entry Ere: \url{https://kaweah.freedombox.rocks/kiwix/wiktionary_en_all_maxi_2023-04/Ere}
\end{itemize}

\section{Religious and Community Resources}
Joshua Project profile of the Ere, Nane people of Manus Province: \url{https://joshuaproject.net/people_groups/11147}

\section{Digital and Audio Resources}
Wikipedia mentions a recording and printed notes held in the Kaipuleohone archive. I did not find the archive's own URL in the results.

\section{Curriculum Development Context}
The reference description says the language is not taught in schools. Tok Pisin is used for wider communication. No orthography, primer or dictionary was found. A curriculum would start from documentation with the community, using the Kaipuleohone holdings as a possible source.
\end{document}